\section{Automated Testing}

The Hitchly project implements automated testing as planned in the VnV Plan Section~3.6 (Automated Testing and Verification Tools) and Section~5 (Unit Test Description). The implementation uses Vitest as the primary testing framework (the plan referenced Jest; the codebase adopted Vitest, which is Jest-compatible). Unit tests live in the backend and shared packages: under \texttt{apps/api/} in \texttt{trpc/routers/\_\_tests\_\_/} and \texttt{services/\_\_tests\_\_/}, and in \texttt{packages/db/tests/}, \texttt{packages/emails/}, and \texttt{packages/utils/}. They cover business logic for authentication and verification, user profile, route and trip, matchmaking, notifications, rating and feedback, safety and reporting, payment and cost estimation, pricing, and the database schema; external dependencies (database, Stripe SDK, Google Maps API) are mocked. Static analysis and formatting are enforced with ESLint and Prettier; the CI pipeline runs \texttt{pnpm lint} and \texttt{pnpm format:check} on every push and pull request.

Continuous integration is implemented in GitHub Actions; the workflow file is \texttt{.github/workflows/ci.yml}. On code changes the workflow runs type-check, lint, format-check, build, and test; the test job uses a Postgres service container and \texttt{pnpm test} to run the Vitest suites across the monorepo. Code coverage is supported by Vitest's integrated reporter (e.g.\ \texttt{vitest --coverage}). The test IDs used elsewhere in this report (including Section~10, Trace to Modules) match the unit test identifiers in the VnV Plan Section~5, providing clear traceability between the plan and the implemented automated testing.
